\documentclass[12pt]{article}
\usepackage[utf8]{inputenc}
\usepackage[T1]{fontenc}
\usepackage{hyperref}
\title{Marimanindji Language Resources for K-12 Curriculum Development}
\author{Compiled by Claude Sonnet 5.5 (claude-sonnet-5-5) for Andr\'as Kornai}
\date{}
\begin{document}
\maketitle

\section{Introduction}
Marimanindji (also Maramanindji; ISO 639-3 zmm, Glottocode mari1417) is a Western Daly language of the Daly River region, Northern Territory. It is dormant or extinct. Perhaps fifteen speakers were reported in 1983, and very few after that. Some sources treat it as a dialect of Marrithiyel. It should not be confused with Maranunggu, a separate language.

\section{Primary Educational Resources}
No primers, textbooks or school programmes were found in search results.

\section{Linguistic and Reference Materials}
\begin{itemize}
\item Wikipedia, Marimanindji language: \url{https://en.wikipedia.org/wiki/Marimanindji_language}
\item Wikipedia, Maramanindji: \url{https://en.wikipedia.org/wiki/Maramanindji}
\item Wikipedia, Western Daly languages: \url{https://en.wikipedia.org/wiki/Western_Daly_languages}
\item Daly Languages site, language page: \url{https://dalylanguages.org/view_language.php?id=16}
\item South Australian Museum archives, Marimanindji language group: \url{https://www.samuseum.sa.gov.au/collection/archives/language_groups/marimanindji}
\end{itemize}
Search results say Tryon's word-list and sketch grammar (1968, 1970, 1974) are the only published information. No URL for them was found.

\section{Religious and Community Resources}
No religious resources were found. The Daly Languages site may point to community contacts.

\section{Digital and Audio Resources}
\begin{itemize}
\item Endangered Languages Project page: \url{https://www.endangeredlanguages.com/lang/1909?hl=en}
\end{itemize}

\section{Curriculum Development Context}
Documentation is minimal. A curriculum would depend on the Tryon word-list, archival holdings and consultation with Daly region communities. Marrithiyel and Brinkin material may be related but should be kept distinct.
\end{document}